\documentclass[11pt]{article}
\usepackage[a4paper,margin=25mm]{geometry}
\usepackage{amsmath,amssymb}
\usepackage[T1]{fontenc}
\usepackage{lmodern}
\usepackage{hyperref}
\setlength{\parindent}{0pt}
\setlength{\parskip}{7pt}
\title{Negative prime-position powers do not define a p-adic number}
\author{SucRunBug}
\date{2 October 2026}
\begin{document}
\maketitle
\textbf{Conjecture 00000000714 -- FALSE AS WRITTEN (DIVERGENT SERIES).}

\textbf{Filed claim addressed.} The series $\Theta=\sum_{n\ge1}p^{-p_n}$ is transcendental in $\mathbb Q_p$.

\section*{Proof}
For any prime $p$ and any prime exponent $p_n$, the $p$-adic norm is
\[|p^{-p_n}|_p=p^{p_n}\ge1.\]
In particular, the terms do not tend to zero in $\mathbb Q_p$. A convergent series in a topological field must have terms tending to zero, because they are successive differences of its convergent partial sums. The displayed series therefore does not converge in $\mathbb Q_p$, and does not define the claimed number $\Theta$.

This is a refutation of the assertion as written, including its presupposition that the displayed sum is a $p$-adic number. Changing the exponents to positive $p_n$ would give a different series; no claim about its transcendence is made here.

\section*{Formalization scope}
Lean uses the actual norm on Padic p, proves the obstruction for every natural-exponent sequence, and proves non-summability for Nat.nth Nat.Prime. This is a convergence obstruction, not a numerical test or a claim about the corrected positive-exponent series.

\textbf{Reproduce.} In \texttt{lean4/}, run \texttt{lake exe cache get},
\texttt{lake build}, then \texttt{lake env lean Check.lean}. Versions:
Lean/Mathlib \texttt{v4.33.0}. Independent finite checks are in
\texttt{reproduce.py}; they are not a substitute for the complete proof.

\textbf{Source.} \url{https://github.com/The-Last-Math-Competition/The-Last-Math-Competition/blob/28a22efac52c4a7abd4bba8a143b1c1a06b7e177/conjectures/00000000714.md}
\end{document}
